\documentclass[11pt]{article}
\usepackage[a4paper,margin=25mm]{geometry}
\usepackage{amsmath,amssymb,amsthm,booktabs,array,hyperref}
\hypersetup{colorlinks=true,urlcolor=blue,linkcolor=blue,citecolor=blue}
\newtheorem{theorem}{Theorem}
\newtheorem{lemma}{Lemma}
\newcommand{\conv}{\operatorname{conv}}
\newcommand{\Vol}{\operatorname{Vol}}
\title{A ten-cube counterexample to the recursive structure conjecture for extremal polyominoes}
\author{Principia Math In Progress}
\date{28 September 2026}
\begin{document}
\maketitle
\begin{abstract}
We give a facet-connected union of ten integer unit cubes with bounding
box side lengths $(4,4,4)$ and convex-hull volume $28$, the maximum possible
for ten cubes. Every coordinate layer contains at most six cubes. Thus no
one-layer subpolyomino can retain both other bounding-box side lengths.
This disproves Conjecture 1 of Kurz, already in dimension three.
\end{abstract}

\section*{Task 1: the source question and its resolution}
In \cite[Conjecture 1, p.~135]{Kurz}, an extremal $d$-dimensional
polyomino is conjectured to contain a subpolyomino $P'$ of thickness one
in some coordinate direction $i$, whose projection retains all other
bounding-box side lengths and is again extremal; the remainder is to
consist of at most two linear arms. The conjecture also asserts that the
auxiliary growth parameters $v_i$ vanish. We disprove the existence of
$P'$ itself, without needing to interpret the arm condition.

For an integer vector $a$ let $Q_a=a+[0,1]^3$. Take the following lower
corners in the displayed growth order, and set $P=\bigcup_{j=0}^9Q_{a_j}$:
\begin{equation}\label{cells}
\begin{array}{c|r@{\,,\,}r@{\,,\,}r|c}
j&\multicolumn{3}{c|}{a_j}&\text{earlier facet neighbor}\\ \hline
0&0&0&0&--\\
1&1&0&0&0\\
2&1&0&-1&1\\
3&1&-1&-1&2\\
4&-1&0&0&0\\
5&1&-2&-1&3\\
6&2&0&0&1\\
7&1&-1&-2&3\\
8&2&0&1&6\\
9&1&-3&-1&5
\end{array}
\end{equation}
Each listed neighbor differs by one standard basis vector up to sign.
Consequently $P$ is a polyomino with ten distinct cubes. Its bounding
box is $[-1,3]\times[-3,1]\times[-2,2]$, of side lengths $(4,4,4)$.
Every added cube increases one bounding-box side length. Thus all the
source's parameters $v_i$ are zero. Equivalently, the source identity
$n=1+\sum_i(l_i-1)+\sum_i v_i$ forces $\sum_i v_i=0$.

\begin{theorem}\label{main}
The polyomino $P$ in \eqref{cells} has the maximum convex-hull volume
among ten-cube three-dimensional polyominoes, but has no subpolyomino
satisfying conditions (i) and (ii) of Kurz's Conjecture 1.
\end{theorem}

\section*{Exact convex-hull volume}
Write $K=\conv\{a_0,\ldots,a_9\}$. Convexity of the unit cube gives
\begin{equation}\label{minkowski}
\conv(P)=K+[0,1]^3.
\end{equation}
The six vertices of $K$ are
\[
\begin{array}{lll}
A=(-1,0,0),&B=(1,0,-1),&C=(1,-1,-2),\\
D=(1,-3,-1),&E=(2,0,0),&F=(2,0,1).
\end{array}
\]
The other four lower corners lie on the segments $AE$ or $BD$.
Here is a complete facet certificate for this six-vertex polytope.
An oriented row $UVW$ contributes $\det(U,V,W)$ to six times the
volume. For the quadrilateral use the fan from its first vertex.
\begin{center}
\begin{tabular}{c c r r}
\toprule
Oriented facet & Outward normal $u$ & $b$ in $u\cdot x=b$
  & Contribution\\
\midrule
$ACD$&$(-5,-2,-4)$&5&5\\
$ADF$&$(-3,-5,9)$&3&3\\
$BCA$&$(-1,2,-2)$&1&1\\
$FEBA$&$(0,1,0)$&0&0\\
$CBE$&$(1,1,-1)$&2&2\\
$DEF$&$(3,-1,0)$&6&6\\
$DCE$&$(5,-1,-2)$&10&10\\
\bottomrule
\end{tabular}
\end{center}
All six vertices satisfy every displayed inequality $u\cdot x\le b$.
Checking the $\binom63=20$ vertex triples gives precisely these seven
supporting planes; thus the table lists every facet. The origin belongs
to $K$ (on $AE$), so pyramids from the origin to the oriented facets
give
\begin{equation}\label{kvol}
\Vol(K)=\frac{5+3+1+0+2+6+10}{6}=\frac92.
\end{equation}
The three coordinate projections have the following cyclic vertex lists
and areas, directly obtained by the shoelace formula. The letters here
denote the corresponding projected points.
\begin{center}
\begin{tabular}{ccc}
\toprule
Coordinate plane & Projected hull & Area\\
\midrule
$xy$&$ADE$&$9/2$\\
$xz$&$ACEF$&$9/2$\\
$yz$&$DCBF$&$9/2$\\
\bottomrule
\end{tabular}
\end{center}
Every coordinate width of $K$ is $3$.

For any three-dimensional convex body $L$, addition of a unit segment
in direction $e_i$ increases its volume by the area of its projection
onto $e_i^\perp$. Indeed, every nonempty parallel line section is an
interval whose length increases by one, and integration over the
projection proves the assertion. Applying this three times, together
with the two-dimensional analogue for projection areas, gives
\begin{align*}
\Vol(K+[0,1]^3)
 &=\Vol(K)+\operatorname{area}(K_{xy})
   +\operatorname{area}(K_{xz})+\operatorname{area}(K_{yz})\\
 &\qquad+\operatorname{width}_x(K)+\operatorname{width}_y(K)
   +\operatorname{width}_z(K)+1\\
 &=\frac92+3\frac92+3+3+3+1=28.
\end{align*}
Kurz's upper bound \cite[Theorem 4]{Kurz} for $n$ cubes in dimension $d$
is
\[
\sum_{I\subseteq\{1,\ldots,d\}}\frac1{|I|!}
  \prod_{i\in I}\left\lfloor\frac{n-2+i}{d}\right\rfloor.
\]
For $d=3$ and $n=10$ the three floor values are $3$, so this equals
$1+9+27/2+27/6=28$. Hence $P$ is globally extremal for ten cubes.

\section*{The missing planar subpolyomino}
\begin{lemma}\label{planar}
A facet-connected planar polyomino with bounding-box side lengths
$r$ and $s$ contains at least $r+s-1$ cells.
\end{lemma}
\begin{proof}
Order its cells by growing a spanning tree of the facet-adjacency graph.
The first cell has side-length sum two. A subsequent cell shares a side
with an earlier cell and can increase that sum by at most one. If there
are $m$ cells, then $r+s\le 2+(m-1)$.
\end{proof}
The layer counts of $P$, indexed by the respective coordinates of its
lower corners, are
\[
\begin{array}{c|c|c}
\text{direction}&\text{layer coordinates}&\text{numbers of cubes}\\\hline
x&-1,0,1,2&1,1,6,2\\
y&-3,-2,-1,0&1,1,2,6\\
z&-2,-1,0,1&1,4,4,1.
\end{array}
\]
In particular, each layer contains at most six cubes. Suppose a
subpolyomino $P'$ had thickness one in direction $i$ and retained both
other side lengths, as conditions (i) and (ii) require. Its cubes would
all lie in one coordinate layer. Projection along $i$ is injective on
these cubes and preserves facet connectivity, so it would give a planar
polyomino with side lengths $4,4$ and at most six cells. Lemma
\ref{planar} requires at least seven. This contradiction proves
Theorem~\ref{main} and disproves the conjecture. \qed

\section*{Verification and limits}
The accompanying standard-library Python verifier uses integer and
rational arithmetic. It checks the growth order, layers, all supporting
planes of $K$ and of the 44 distinct cube corners, oriented facet
volumes, edge incidences, Euler's formula, projection areas, and the
source bound. The direct cube-corner hull has 20 vertices, 37 edges and
19 facets, and volume $28$, independently confirming the extrusion
calculation. It also exhausts all 1,023 nonempty subsets of the ten cubes
for an eligible planar subpolyomino. Numerical search was used only to
find the example. No smallest-counterexample claim is made. This is an
author-agent proof and audit, without independent peer review.

\section*{Proposed Tasks 2--5}
\begin{enumerate}
\setcounter{enumi}{1}
\item Determine the smallest number of cubes in a three-dimensional
counterexample to Kurz's recursive structure conjecture, with an exact
classification or exclusion certificate for every smaller size.
\item Classify all ten-cube convex-hull maximizers up to translation,
coordinate permutation and reflection, and determine which admit a
subpolyomino satisfying the conjectured recursion.
\item Construct infinitely many sizes for which a three-dimensional
convex-hull maximizer has no coordinate layer containing a connected
subpolyomino with both full remaining side lengths.
\item Determine for each dimension $d\ge4$ whether such a failure of the
full-side-length layer condition occurs among convex-hull maximizers,
giving constructions or a structural obstruction.
\end{enumerate}

\begin{thebibliography}{9}
\bibitem{Kurz} S.~Kurz, \emph{Convex hulls of polyominoes},
Beitr\"age zur Algebra und Geometrie \textbf{49} (2008), no.~1, 125--136.
\href{https://www.maths.tcd.ie/EMIS/journals/BAG/vol.49/no.1/b49h1kur.pdf}{Published article};
\href{https://arxiv.org/abs/math/0702786}{arXiv:math/0702786v1} (2007).
\bibitem{MathDB} \href{https://mathdb.com/p/380444}{MathDB 380444},
\emph{Recursive structure conjecture for extremal polyominoes}.
Source and public-claim audit dated 28 September 2026; details accompany
the proof package.
\end{thebibliography}
\end{document}
